%=== tech/fonctions/lectgraphiques_01.tex || Fonctions || Graphique

\hfill\begin{GraphiqueTikz}[x=0.375cm,y=0.375cm,Xmin=-3,Xmax=4,Xgrilles=1,Ymin=-3,Ymax=4,Ygrilles=1]
	\TracerAxesGrilles[Origine,Police=\tiny,Origine]{-2,1,3}{1}
	\TracerCourbe[Couleur=blue,Debut=-2,Fin=3]< f >{x}
	\TracerCourbe[Couleur=red,Debut=-2,Fin=1]< g >{4/3*x^(2)+x-4/3}
	\TracerCourbe[Couleur=red,Debut=1,Fin=3]< h >{-(x-2)^2+2}
	\MarquerPts*[Style=+]{(-2,-2),(-1,-1),(1,1),(3,3),(3,1),(-2,2)}
	\PlacerTexte[Couleur=blue,Police=\footnotesize,Position=above left]{(3,2.5)}{$\mathcal{C}^{\prime}$}
	\PlacerTexte[Couleur=red,Police=\footnotesize,Position=below right]{(-2,2)}{$\mathcal{C}$}
\end{GraphiqueTikz}\hfill\null

\smallskip

Dans la figure ci-dessus, les courbes $\mathcal{C}$ et $\mathcal{C}^{\prime}$ représentent respectivement les fonctions $f$ et $g$.

L'ensemble des solutions de l'inéquation $f(x) \leqslant g(x)$ est :

\smallskip

\eamreponsesautom[2]{$[-2 ;-1]$}{$[1 ; 2]$}{$[-2 ;-1] \cup[1 ; 2]$}{$[-2 ;-1] \cap[1 ; 2]$}
